\section{Young evolution with constant initial data}
\label{S:constant-phase-young}

We use the fixed $R,T,E,\gamma,\theta,\beta$ of~\eqref{E:driver-parameters}, and
choose the allowed even mask $m = \zeta_R$ with $0 \le m \le 1$. All constants
in this section can depend on these fixed data, the model and the auxiliary
physical cutoffs. There is no uniform estimate as $R$ or $T$ grows.

\subsection{The distributional phase and its free evolution}

Define the real Hilbert space
\begin{equation}\label{E:constant-phase-space}
  F=E\oplus\R^2,\qquad
  \|(e,c_1,c_2)\|_F^2=\|e\|_E^2+c_1^2+c_2^2,\qquad
  e_0=(0,1,0).
\end{equation}
Its distributional interpretation is $(e,c_1,c_2)\mapsto(e_1+c_1,e_2+c_2)$. Let
$i:E \to F$ denote $e\mapsto(e,0,0)$ and let $\pi_C:F \to \R^2$ be the
constant-coordinate projection.

\begin{lemma}[Constants and finite propagation]\label{L:constant-phase-free}
  The distributional interpretation of $F$ is continuous and injective. The
  space is complete and separable, and $iE$ is closed and proper. For every real
  $\psi \in C_c^\infty(\R^d)$,
  \begin{equation}\label{E:constant-phase-cutoff}
    L_\psi(e,c_1,c_2)=\psi e+(c_1\psi,c_2\psi)
  \end{equation}
  is bounded from $F$ to $E$. The free wave group on $F$ is
  \begin{equation}\label{E:constant-phase-free}
    S_F(t)(e,c_1,c_2)=(S(t)e,c_1+tc_2,c_2).
  \end{equation}
  It is strongly continuous and bounded on finite horizons and has finite
  propagation for the total distributional data, including local cancellation
  between $e$ and the constants.
\end{lemma}

\begin{proof}
  A nonzero constant has Fourier transform a nonzero multiple of $\delta_0$.
  Every $H^\sigma$ Fourier transform is locally $L^2$ near zero, since its
  weight is bounded above and below there. A locally $L^2$ function $f$ cannot
  represent $\delta_0$: testing against $\varphi(\xi/\varepsilon)$, with
  $\varphi(0) = 1$ and $\varphi$ compactly supported in $B_1$, bounds its
  integral by
  \[
    \|f\|_{L^2(B_\varepsilon)}
    \varepsilon^{d/2}\|\varphi\|_2\longrightarrow0,
  \]
  whereas $\delta_0$ gives one. Thus, constants lie in neither $H^{1/2}$ nor
  $H^{-1/2}$; the interpretation is injective. Sobolev duality and the two test
  integrals prove continuity into distributions, uniformly on bounded families
  of compact tests. Completeness, separability and the closed proper inclusion
  follow from the Hilbert direct sum.

  Multiplication by $\psi$ is bounded on both Sobolev entries. Indeed, the
  Peetre inequalities give $\langle\xi\rangle^\sigma\langle\eta\rangle^{-\sigma}
  \le 2^{|\sigma|/2}\langle\xi-\eta\rangle^{|\sigma|}$. In the unitary Fourier
  convention its weighted kernel is therefore bounded by convolution with the
  $L^1$ function $2^{|\sigma|/2}c_d|\widehat{\psi}(q)|\langle
  q\rangle^{|\sigma|}$, where $c_d = (2\pi)^{-d/2}$. The $L^1*L^2$ bound applies
  at $\sigma = \pm1/2$; the two constant columns are also in $E$. This proves
  boundedness of~\eqref{E:constant-phase-cutoff}.

  The finite-dimensional block in~\eqref{E:constant-phase-free} and the strong,
  finite-horizon bound for $S(t)$ from Lemma~\ref{L:driver-gain-localization}
  prove the group assertions. To verify its physical meaning, for $t > 0$
  consider the matrix of compactly supported distributions
  \[
    \begin{pmatrix}
      \partial_tG_d(t)&G_d(t)\\
      \Delta G_d(t)&\partial_tG_d(t)
    \end{pmatrix}.
  \]
  The source Green convention gives $\widehat{G}_d(t)(\xi) =
  c_d\sin(t|\xi|)/|\xi|$; hence its convolution multiplier is the sine quotient,
  and the differentiated multipliers are $\cos(t|\xi|)$ and
  $-|\xi|\sin(t|\xi|)$. The matrix agrees with $S(t)$ on Schwartz phases.
  Convolution by a compact distribution is continuous on tempered distributions:
  its transpose convolves a Schwartz test with the reflected compact
  distribution, continuously in all Schwartz seminorms. Schwartz approximation
  in $E$ extends the identity to $E$. The masses of the four entries are
  $1,t,0,1$, so constants give exactly~\eqref{E:constant-phase-free}. Time
  reversal handles $t < 0$. Each entry is supported in $B_{|t|}$, also for the
  three-dimensional sphere kernel. If the total initial distribution vanishes on
  an enlarged region, each translated kernel tests zero there, so the evolved
  total distribution vanishes on the smaller cone. This argument permits
  cancellation between its Sobolev and constant parts. This proves
  Lemma~\ref{L:constant-phase-free}.
\end{proof}

\subsection{The two columns for constant initial data}

Choose real $\chi \in C_c^\infty$, equal to one on $B_{2R+T+1}$ and supported in
$B_{2R+T+2}$, and choose real compact smooth $\psi = 1$ on a neighborhood of
$\operatorname{supp}\chi$. Write $M_E = M_R$ and $M_E^\rho = M_{Y_\rho}$ for the
paths already constructed, with the same original deterministic coefficients.

\begin{proposition}[Extended operator path]\label{P:constant-phase-driver}
  As identities of continuous operator paths,
  \begin{equation}\label{E:constant-phase-localization}
    D_\chi M_E=M_E=M_ED_\chi,
  \end{equation}
  and the same holds for each fixed $\rho$. Define
  \begin{equation}\label{E:constant-phase-driver}
    A_E=M_EL_\psi,\qquad A=iA_E,\qquad
    A_E^\rho=M_E^\rho L_\psi,\qquad A^\rho=iA_E^\rho.
  \end{equation}
  These are independent of $\psi$, have compact values, and satisfy $Ai = iM_E$,
  $\pi_CA = 0$ and their source counterparts. In every finite $L^pC^\theta$,
  $A^\rho \to A$ at arbitrary relative source-width rates. Their norms have
  uniform $\sqrt p$ moment bounds and a common finite square-exponential moment.
  The original $A$ is measurable with respect to the given $W$.
\end{proposition}

\begin{proof}
  Lemma~\ref{L:driver-gain-localization} gives both localizations for every
  deterministic $A_j(0,t)$. The original finite sums retain them.
  Proposition~\ref{P:source-gaussian-transfer} constructs the source finite sums
  with exactly these coefficients and correlated $Y_\rho(f_j)$. Their $C^\theta$
  limits preserve both identities. A uniformly convergent subsequence, or
  rational times and continuous paths,
  gives~\eqref{E:constant-phase-localization} for all $t$ on one event. Use a
  separate event for each fixed $\rho$.

  For two choices $\psi,\psi'$, $(L_\psi-L_{\psi'})z \in E$ and
  $D_\chi(L_\psi-L_{\psi'})z = 0$. The right localization kills this difference.
  For $z = ie$, $D_\chi D_\psi e = D_\chi e$ and hence $Ai = iM_E$. The constant
  columns are precisely $iM_E(\psi,0)$ and $iM_E(0,\psi)$; they are retained.
  Bounded composition preserves compactness and transfers
  \eqref{E:source-driver-limit} and the Gaussian norm bounds
  to~\eqref{E:constant-phase-driver}. If $\|X\|_{L^p} \le C\sqrt p$, then $\E
  X^{2n} \le C^{2n}(2n)^n$. Since $n! \ge (n/e)^n$, the exponential series gives
  $\E e^{cX^2} < \infty$ uniformly when $2ecC^2 < 1$. Finally, $A$ is a
  deterministic bounded transform of the original $W$-measurable operator path.
  This proves Proposition~\ref{P:constant-phase-driver}.
\end{proof}

\begin{lemma}[Ordinary response on all of $F$]\label{L:constant-source-response}
  For fixed $\rho$ put $b(r) = mY_\rho(r)$ and define
  \[
    B_F(b)(e,c_1,c_2)=(0,b(e_1+c_1))\in E.
  \]
  Then
  \begin{equation}\label{E:constant-source-response}
    A_E^\rho(t)z=\int_0^tS(-r)B_F(b(r))S_F(r)z\ud r
  \end{equation}
  as a statewise strong integral. Its primitive is Lipschitz in
  $\mathcal{B}(F,E)$ and the identity holds for all states and times on one
  event for this fixed $\rho$.
\end{lemma}

\begin{proof}
  The source regularity gives one compact spatial support for $b$, continuity in
  $L^\infty$ and $L^2$, and bounded spatial derivatives on the horizon. Since
  $e_1 \in L^2$ and $\|\cdot\|_{H^{-1/2}} \le \|\cdot\|_2$,
  \[
    \|B_F(b)\|_{\mathcal{B}(F,E)}
    \le\|b\|_\infty+\|b\|_2.
  \]
  Strong continuity and finite-horizon bounds of both free groups make the
  integrand strongly continuous on each state and uniformly bounded as an
  operator. Integrating its bound proves operator-norm Lipschitz increments
  without assuming a Banach-valued operator derivative.

  The total initial distributions of $z$ and $iL_\psi z$ agree on the entire
  enlarged ball where $\psi = 1$. Lemma~\ref{L:constant-phase-free} gives
  equality of the free displacements near $\operatorname{supp}m$ for $0 \le r
  \le T$. Multiplication by $b(r)$ kills the difference. Apply
  Proposition~\ref{P:source-ordinary-response} to the $E$ state $L_\psi z$ to
  obtain~\eqref{E:constant-source-response}. Separability and rational times,
  followed by boundedness and continuity, give its path identity. This proves
  the constant columns directly; no density of $E$ in $F$ is used. This proves
  Lemma~\ref{L:constant-source-response}.
\end{proof}

\subsection{The exact linear Young equation}

We use Young's integration theory~\cite{young:36:inequality} in its
Banach-space form. The estimate is stated by
Friz and Hairer~\cite{friz.hairer:20:course}*{Sections~1.4--1.5, 4.1, (4.3), pp.~13, 61--62}.
For $D,V$ with exponent $\theta > 1/2$ and bounded bilinear multiplication, left
composition gives
\begin{equation}\label{E:constant-young-estimate}
  \left\|\int_s^t\ud D(r)V(r)-D(s,t)V(s)\right\|
  \le C_\theta[D]_\theta[V]_\theta|t-s|^{2\theta}.
\end{equation}
The Banach spaces here are $\mathcal{B}(F)$, and multiplication is bounded with
constant one. We derive the linear solution and stability estimates below rather
than invoke a nonlinear bounded-vector-field theorem.

\begin{proposition}[Linear Young evolution]\label{P:constant-young-linear}
  For every $D \in C^\theta([0,T];\mathcal{B}(F))$ with $D(0) = 0$ there is a
  unique $C^\theta$ solution, in the terminal variable, of
  \begin{equation}\label{E:constant-young-equation}
    U_D(t,s)=I+\int_s^t\ud D(r)U_D(r,s).
  \end{equation}
  It has the composition law, invertible values and joint operator-norm
  continuity on $0 \le s \le t \le T$. The map $D\mapsto U_D$ is locally
  Lipschitz into the uniform time-triangle norm, and
  \begin{equation}\label{E:constant-young-growth}
    \sup_{s\le t\le T}\|U_D(t,s)\|
    \le C_T\exp\big(C_T(1+\|D\|_{C^\theta})^{1/\theta}\big).
  \end{equation}
\end{proposition}

\begin{proof}
  On an interval $[a,a+\ell]$ use the complete scaled norm $N_\ell(V) =
  \sup\|V\|+\ell^\theta[V]_\theta$ and put $I_D(V)(t) = \int_a^t\ud D\,V$.
  Applying~\eqref{E:constant-young-estimate} to every subinterval gives
  \begin{align*}
    \sup\|I_D(V)\|
    &\le[D]_\theta\ell^\theta\sup\|V\|
    +C_\theta[D]_\theta\ell^{2\theta}[V]_\theta,\\
    [I_D(V)]_\theta
    &\le[D]_\theta\sup\|V\|
    +C_\theta[D]_\theta\ell^\theta[V]_\theta.
  \end{align*}
  Thus, with $c_\theta = 2\max(1,C_\theta)$,
  \begin{equation}\label{E:constant-young-local}
    N_\ell(I_D(V))
    \le c_\theta[D]_\theta\ell^\theta N_\ell(V).
  \end{equation}
  Take $K \ge 1$, $[D]_\theta \le K$, and choose $\ell$ with $c_\theta
  K\ell^\theta \le 1/4$. The map $V\mapsto V_a+I_D(V)$ is a contraction on the
  complete affine space of paths starting at $V_a$. Its solution satisfies
  $N_\ell(V) \le (4/3)\|V_a\| \le 2\|V_a\|$.

  To compare $D,D'$, require both seminorms at most $K$ and write $\delta =
  \|D-D'\|_{C^\theta}$, $W = V-V'$. Subtracting gives $W =
  W_a+I_D(W)+I_{D-D'}(V')$. Absorbing the factor $1/4$ yields
  \begin{equation}\label{E:constant-young-difference}
    N_\ell(W)\le2\|W_a\|
    +2c_\theta\ell^\theta\delta N_\ell(V').
  \end{equation}
  Cover $[s,T]$ by at most $n_{K,T}$ such pieces. If the second initial norm is
  $B$, its norm on piece $j$, indexed from zero, is at most $2^{j+1}B$. With
  $E_j$ the difference at its left endpoint, induction gives
  \[
    E_{j+1}\le2E_j+2^{j+2}c_\theta\ell^\theta\delta B,\qquad
    E_j\le2^jE_0+j2^{j+1}c_\theta\ell^\theta\delta B.
  \]
  Equation~\eqref{E:constant-young-difference} gives the same bound inside each
  piece; a shorter final piece only reduces the coefficient. For evolutions,
  $E_0 = 0$ and $B = 1$. The number of pieces is uniform in $s,t$, proving local
  Lipschitz dependence in the full triangle norm.

  Now choose $K = 1+\|D\|_{C^\theta}$. At most
  $C_T(1+\|D\|_{C^\theta}^{1/\theta})$ pieces suffice. Each multiplies the norm
  by at most two, giving \eqref{E:constant-young-growth}. Uniqueness patches the
  solutions; right multiplication by the preceding evolution gives the
  composition law. On a small piece~\eqref{E:constant-young-local} gives
  $\|U_D-I\|_\infty \le 1/3$, so a Neumann series proves local invertibility,
  then products give global invertibility. Finally, $U_D(t,s) =
  U_D(t,0)U_D(s,0)^{-1}$ proves joint norm continuity, since $U_D(\cdot,0)$ is
  continuous and inversion is continuous on invertible operators. This proves
  Proposition~\ref{P:constant-young-linear}.
\end{proof}

\begin{proposition}[Gaussian Young and physical flows]
  \label{P:constant-gaussian-flow}
  Apply Proposition~\ref{P:constant-young-linear} to $A,A^\rho$. The resulting
  flows are strongly measurable in uniform triangle norm, have every finite norm
  moment and satisfy
  \begin{equation}\label{E:constant-flow-limit}
    U_{A^\rho}\longrightarrow U_A
    \quad\hbox{in every finite }L^p
    \left(\sup_{0\le s\le t\le T}\|\cdot\|_{\mathcal{B}(F)}\right).
  \end{equation}
  They preserve the interaction constant coordinates:
  \begin{equation}\label{E:constant-flow-range}
    \pi_CU_A(t,s)=\pi_C,\qquad
    (U_A(t,s)-I)F\subset iE,
  \end{equation}
  and likewise for the source. The physical flow is
  \begin{equation}\label{E:constant-physical-flow}
    P_A(t,s)=S_F(t)U_A(t,s)S_F(-s).
  \end{equation}
  It has composition and strong continuity on states, physical constant block
  $(c_1,c_2)\mapsto(c_1+(t-s)c_2,c_2)$, all finite moments and the analogous
  uniform $L^p$ source convergence. Its difference from the free flow has range
  in $iE$.
\end{proposition}

\begin{proof}
  The original paths lie in separable closed path spans by their Gaussian
  construction. Bounded composition in~\eqref{E:constant-phase-driver} and the
  continuous Young map preserve strong measurability into
  $C(\Delta_T;\mathcal{B}(F))$, where $\Delta_T = \{(t,s):0\le s\le t\le T\}$
  has the uniform norm. Indeed, the continuous image of the separable space of
  operator paths has separable closure in this triangle space. Applying $\pi_C$
  to Young Riemann sums gives the first identity
  in~\eqref{E:constant-flow-range}. Each increment applied to any state lies in
  $iE$; closedness of this subspace passes the property through sums and limits,
  giving the second identity.

  Write $X_\rho = \|A^\rho\|_{C^\theta}$ and $X_0 = \|A\|_{C^\theta}$.
  Proposition~\ref{P:constant-phase-driver} supplies $c_0 > 0$ and $H_0 <
  \infty$ such that
  \[
    \sup_\rho\E e^{c_0X_\rho^2}\le H_0,\qquad
    \E e^{c_0X_0^2}\le H_0.
  \]
  For $K \ge 1$, let $L_K$ be the Lipschitz constant of
  Proposition~\ref{P:constant-young-linear} on the ball of operator paths of
  radius $K$, in the whole time-triangle norm. On $\{X_\rho\le K,\ X_0\le K\}$
  it gives
  \[
    \|U_{A^\rho}-U_A\|_\infty
    \le L_K\|A^\rho-A\|_{C^\theta}.
  \]
  Consequently, for $\varepsilon > 0$ and any finite $q \ge 1$, Markov and the
  union bound give
  \[
    \mathbb{P}(\|U_{A^\rho}-U_A\|_\infty>\varepsilon)
    \le2H_0e^{-c_0K^2}
    +(L_K/\varepsilon)^q\E\|A^\rho-A\|_{C^\theta}^q.
  \]
  At fixed $K$ the second term tends to zero at arbitrary relative source-width
  rates. Only then let $K \to \infty$ to obtain convergence in probability. No
  independence of the original, source and residual operator paths is required.

  Put $\nu = 1/\theta < 2$. For each finite $r \ge 1$ and $c > 0$, there is
  $C_{r,c} < \infty$ such that
  \[
    rC_T(1+x)^\nu\le cx^2+C_{r,c},\qquad x\ge0.
  \]
  This follows because $(1+x)^\nu/x^2 \to 0$ as $x \to \infty$, while the
  remainder is bounded on compact intervals. Use $c \le c_0$
  and~\eqref{E:constant-young-growth} to obtain
  \[
    \sup_\rho\E\|U_{A^\rho}\|_\infty^r<\infty,\qquad
    \E\|U_A\|_\infty^r<\infty.
  \]
  In particular, for fixed finite $p \ge 1$ and $V_\rho =
  \|U_{A^\rho}-U_A\|_\infty$, the $2p$ moments of $V_\rho$ are uniformly
  bounded. For every $\varepsilon > 0$,
  \[
    \E V_\rho^p
    \le\varepsilon^p+
    (\E V_\rho^{2p})^{1/2}
    \mathbb{P}(V_\rho>\varepsilon)^{1/2}.
  \]
  The second term tends to zero by convergence in probability; then let
  $\varepsilon\downarrow0$. This proves \eqref{E:constant-flow-limit}. The
  square-exponential bound used here is an operator path bound; no
  square-exponential flow bound is needed.

  The free group is strongly continuous and bounded on the horizon. Define
  $\mathcal{B}_\infty(\Delta_T;\mathcal{B}(F))$ to be the Banach space of all
  bounded operator families on $\Delta_T$. The deterministic map
  \[
    \mathcal{L}(V)(t,s)=S_F(t)V(t,s)S_F(-s)
  \]
  is bounded linear from $C(\Delta_T;\mathcal{B}(F))$ into this bounded-family
  space, with norm at most $C_T^2$, where $C_T \ge \sup_{|r|\le T}\|S_F(r)\|$.
  Hence, it preserves strong measurability on the separable ranges and gives, in
  particular,
  \[
    \|P_{A^\rho}-P_A\|_\infty
    \le C_T^2\|U_{A^\rho}-U_A\|_\infty.
  \]
  All finite moments and uniform $L^p$ convergence therefore transfer to the
  physical flows. This does not place those flows in the norm-continuous
  triangle space or assert operator-norm continuity of the free group. For each
  fixed state, joint continuity follows by splitting a physical-flow difference
  into the norm-continuous $U$ difference and two strong free-group differences,
  using the uniform group bound. Group cancellation gives composition. The
  finite constant block $J_t(c_1,c_2) = (c_1+tc_2,c_2)$ satisfies $J_tJ_{-s} =
  J_{t-s}$, which gives the asserted physical constant coordinates from
  $\pi_CU_A = \pi_C$. Finally,
  \[
    P_A(t,s)-S_F(t-s)
    =S_F(t)(U_A(t,s)-I)S_F(-s)
  \]
  has range in $iE$, since $S_F(t)iE = iE$. This proves
  Proposition~\ref{P:constant-gaussian-flow}.
\end{proof}

\subsection{Ordinary source solutions and the original tested observable}

\begin{lemma}[Identification with the ordinary source equation]
  \label{L:constant-source-solution}
  For fixed $\rho$, the displacement distribution of $P_{A^\rho}(t,0)e_0$ is the
  locally continuous ordinary source solution $u_\rho^m(t,\cdot)$ from
  Proposition~\ref{P:masked-source-moments}.
\end{lemma}

\begin{proof}
  In~\eqref{E:constant-source-response} let $C(r) = S(-r)B_F(b(r))S_F(r)$. The
  primitive is $A_E^\rho$. For left endpoints $r_j = t_j$, the Young sums
  applied to a state $v$ are
  \[
    \sum_j\int_{t_j}^{t_{j+1}}iC(r)U_{A^\rho}(r_j,s)v\ud r.
  \]
  Uniform continuity of $U$ and the uniform bound for $C$ make these converge to
  the strong integral with $U(r,s)v$. Thus, the Young equation is precisely the
  ordinary interaction integral on states. Multiplying by free groups gives the
  Duhamel equation for displacement and velocity. Its unique $C([0,T];F)$
  solution follows by the bounded-coefficient Picard argument with the strongly
  continuous, bounded free group. No operator-norm derivative of $S_F$ is used.

  More explicitly, put $W_0(t) = e_0$ and, for $n \ge 1$,
  \[
    W_n(t)=\int_0^t S_F(t-r)iB_F(b(r))W_{n-1}(r)\ud r.
  \]
  If $B_T = \sup_r\|B_F(b(r))\|_{\mathcal{B}(F,E)}$, induction on the time
  simplex gives $\sup_t\|W_n(t)\|_F \le C_T^{n+1}B_T^nT^n/n!$. We take $C_T \ge
  \max(1,\sup_{|t|\le T}\|S_F(t)\|)$, which also supplies the bound at order
  zero. The series converges in $C([0,T];F)$. The Duhamel map there has norm at
  most $C_TB_TT$, so its sum solves the same integral equation and equals the
  ordinary phase by uniqueness.

  Let $X_0(t,x) = 1$ and $X_n(t) = \int_0^tG_d(t-r)*(b(r)X_{n-1}(r))\ud r$. At
  fixed $n$ these are continuous in time and smooth in space. Derivatives fall
  on the smooth compact factor $bX_{n-1}$; Green mass $t-r$ bounds the
  integrals, while the unit-time measure representation gives time continuity.
  This is induction using only finitely many bounded derivatives for each fixed
  derivative order. At $n = 0$ the phase displacement is $X_0$. If the equality
  holds at $n-1$, multiplication by smooth $b$ identifies the forcing with
  $(0,f_r)$, where $f_r = b(r)X_{n-1}(r)$ is smooth and compactly supported. For
  a compact smooth test $\varphi$, continuous distributional testing of the
  phase integral yields
  \[
    \langle[W_n(t)]_1,\varphi\rangle
    =\int_0^t\int f_r(y)
    \int G_d(t-r,\ud z)\varphi(y+z)\ud y\ud r
    =\langle X_n(t),\varphi\rangle.
  \]
  Absolute Fubini follows from $\|\varphi\|_\infty\|f_r\|_1(t-r)$, integrable on
  the horizon because $f_r$ has one compact support and continuous $L^1$ norm.
  This also covers sphere measures in $d = 3$. Iterating the scalar recursion
  gives exactly $X_{\rho,n}^m$, with chronological order and all physical mask
  factors. The scalar compact-cone factorial bound gives locally uniform
  convergence to $u_\rho^m$, whereas the phase series converges in $F$ and hence
  in distributions. Their displacement limits agree. This proves
  Lemma~\ref{L:constant-source-solution}.
\end{proof}

\begin{proposition}[Original average and original point moment]
  \label{P:constant-raw-average}
  Let real $\phi \in C_c^\infty(\R^d)$ be nonnegative with integral one. The
  bounded functional and original random variable
  \begin{equation}\label{E:constant-raw-average}
    O_\phi(e,c_1,c_2)=\langle e_1,\phi\rangle+c_1,\qquad
    Z_m(t,\phi)=O_\phi P_A(t,0)e_0
  \end{equation}
  are well defined, and $Z_m$ is $W$-measurable. At arbitrary relative
  source-width rates, for every finite $p$,
  \begin{equation}\label{E:constant-tested-limit}
    \int\phi(x)u_\rho^m(t,x)\ud x
    \longrightarrow Z_m(t,\phi)\quad\text{in }L^p.
  \end{equation}
  In particular,
  \begin{equation}\label{E:constant-original-point-transfer}
    \E|Z_m(t,\phi)|^2\le\E|u(t,0)|^2.
  \end{equation}
\end{proposition}

\begin{proof}
  Sobolev duality gives $\|O_\phi\| \le (\|\phi\|_{H^{-1/2}}^2+1)^{1/2}$. The
  original operator path is $W$-measurable, and the Young and deterministic maps
  above are continuous or bounded measurable on their separable ranges. Thus,
  $Z_m$ is $W$-measurable and does not depend on the auxiliary $J$ used in
  source realizations. At each fixed $\rho$,
  Lemma~\ref{L:constant-source-solution} identifies the displacement as a
  distribution with the ordinary source solution. Testing by this same $\phi$,
  with $\int\phi = 1$, therefore gives the almost sure identity
  \[
    O_\phi P_{A^\rho}(t,0)e_0
    =C_\phi:=\int\phi(x)u_\rho^m(t,x)\ud x.
  \]
  This uses the displacement pairing and its constant displacement coordinate
  in~\eqref{E:constant-raw-average}; it takes no point trace of an $H^{1/2}$
  state. All identities here are at one fixed $\rho$, so no uncountable
  intersection of source-version events is used.

  We verify that $C_\phi$ equals the Bochner $L^2$ average $B_\phi$ from
  Proposition~\ref{P:masked-source-moments}. That proposition gives the
  spatially uniform bound $\|u_\rho^m(t,x)\|_2 \le Q_t$, with $Q_t$ also
  independent of source widths. Scalar Fubini gives
  \[
    \E\int|\phi(x)u_\rho^m(t,x)|\ud x
    \le Q_t\|\phi\|_1<\infty.
  \]
  Thus, the ordinary spatial integral exists almost surely. Weighted
  Cauchy--Schwarz and scalar Fubini also give
  \[
    \E|C_\phi|^2
    \le\|\phi\|_1\int|\phi(x)|\E|u_\rho^m(t,x)|^2\ud x
    \le Q_t^2\|\phi\|_1^2.
  \]
  For every bounded real random variable $\Xi$ on the fixed source space,
  $V\mapsto\E[\Xi V]$ is a continuous functional on $L^2$. It therefore commutes
  with the Bochner integral, while scalar Fubini is justified by $\|\Xi\|_\infty
  Q_t\|\phi\|_1$. Consequently,
  \[
    \E[\Xi B_\phi]
    =\int\phi(x)\E[\Xi u_\rho^m(t,x)]\ud x
    =\E[\Xi C_\phi].
  \]
  Both averages are real $L^2$ variables. Taking $\Xi =
  \operatorname{sign}(B_\phi-C_\phi)$ yields $\E|B_\phi-C_\phi| = 0$, so the two
  averages agree almost surely.

  Proposition~\ref{P:constant-gaussian-flow} now gives, for every finite $p$,
  \[
    \|B_\phi-Z_m(t,\phi)\|_{L^p}
    \le\|O_\phi\|\|e_0\|
    \big\|\|P_{A^\rho}-P_A\|_\infty\big\|_{L^p}
    \longrightarrow0.
  \]
  This is~\eqref{E:constant-tested-limit}, including its $L^2$ case, for the
  masked Bochner average, with arbitrary relative source-width rates. It
  supplies the conditional convergence premise of
  Proposition~\ref{P:masked-canonical-transfer} at the same fixed mask, horizon,
  original covariance and observable. Since $\phi \ge 0$ and $\|\phi\|_1 = 1$,
  that proposition gives \eqref{E:constant-original-point-transfer}. Zero
  covariance gives the free constant solution and equality. All source limits
  are at fixed mask and horizon, before time asymptotics. This proves
  Proposition~\ref{P:constant-raw-average}.
\end{proof}

\begin{proposition}[Compact deterministic base]\label{P:constant-compact-base}
  Let $h \in C([0,T];C_c^\infty(K))$ be real for one fixed compact $K$, with all
  spatial derivatives continuous on that support. The statewise primitive
  \[
    A_h(t)=\int_0^t iS(-r)B_F(h(r))S_F(r)\ud r
  \]
  is Lipschitz in $\mathcal{B}(F)$, takes values with range in $iE$, and has a
  bound $C_T\sup_r(\|h(r)\|_\infty+\|h(r)\|_2)$. The flow conclusions above,
  including moments, source convergence and ordinary-source identification, hold
  for $A_h+gA$ and $A_h+gA^\rho$, uniformly on bounded displayed deterministic
  bounds and bounded real $g$. No Cameron-shift identity is asserted here.
\end{proposition}

\begin{proof}
  The estimate for $B_F$ and the strong free groups prove the primitive and its
  Lipschitz bound exactly as in Lemma~\ref{L:constant-source-response}. If
  $\|A_h\|_{C^\theta} \le H$, $|g| \le G$ and $X = \|A\|_{C^\theta}$, then
  $\|A_h+gA\|_{C^\theta} \le H+GX$. If $\E e^{cX^2}$ is bounded, choose $c' \le
  c/(2(1+G^2))$ to obtain
  \[
    \E e^{c'\|A_h+gA\|_{C^\theta}^2}
    \le e^{2c'H^2}\E e^{2c'G^2X^2}<\infty.
  \]
  Alternatively $pC(1+H+Gx)^{1/\theta} \le cx^2+C_{p,H,G,c}$ transfers the Young
  growth estimate directly to all finite moments. Source operator-path
  convergence, local Young stability and uniform integrability give the same
  flow convergence. The state-integral and Picard-to-chain arguments apply to
  $h+gmY_\rho$, whose spatial support lies in one fixed compact union. These are
  individual estimates and versions, without an uncountable common-parameter
  event. Only $h = 0$, $g = 1$ enters
  \eqref{E:constant-original-point-transfer}. This proves
  Proposition~\ref{P:constant-compact-base}.
\end{proof}
